\chapter{洛朗级数与留数定理}
\label{ch:b3:residues}

全纯函数在某点附近，且该点是它
\emph{没有定义}的点时，会怎样？
答案是一个完整的三分法——可去奇点、极点与本质奇点——由
双边幂级数，即洛朗展开，直接读出。该展开的一个系数
即\emph{留数}，控制着每一条
绕奇点的围道积分：留数定理把
困难的定积分化为有限的代数运算，计算
函数的零点个数（辐角原理、儒歇定理），并证明开
映射定理。我们先把柯西定理从
星形区域推广到它最终的、不依赖同调的形式——
狄克逊的巧妙论证——从而使所有
相对于补集环绕数为零的围道均可使用。

\section{整体柯西定理}

\emph{轮换}\index{轮换}$\Gamma$ 是闭路径
$\gamma_1, \dots, \gamma_m$ 的有限形式和；
沿 $\Gamma$ 的积分与指标分别为对应各项之和，且
$\operatorname{im}\Gamma = \bigcup\operatorname{im}\gamma_j$。

\begin{theorem}[柯西，整体形式]
\label{thm:b3:residues:globalcauchy}
设 $\Omega \subseteq \C$ 为开集，$f \in \mathcal H(\Omega)$，
且 $\Gamma$ 为 $\Omega$ 中的一个轮换，满足
\[
\operatorname{Ind}_\Gamma(w) = 0
\qquad\text{对每个 } w \notin \Omega .
\]
则对所有 $z \in \Omega\setminus\operatorname{im}\Gamma$，
\[
\frac1{2\iu\pi}\int_\Gamma\frac{f(w)}{w - z}\,\dd w =
\operatorname{Ind}_\Gamma(z)\,f(z),
\qquad\text{且}\qquad
\int_\Gamma f(w)\,\dd w = 0 .
\]
\index{柯西定理（整体）}
\end{theorem}

\begin{proof}[狄克逊证明]
定义 $g \colon \Omega\times\Omega \to \C$ 为
\[
g(z, w) = \begin{cases}
\dfrac{f(w) - f(z)}{w - z} & w \neq z,\\[4pt]
f'(z) & w = z .
\end{cases}
\]
\emph{$g$ 连续}：对角线外显然；在
对角线上一点 $(a, a)$ 附近，将 $f$ 在 $a$ 处展开为幂级数
（\cref{thm:b3:holomorphic:analytic}）：$f(w) - f(z) =
\sum_{n\geq1}c_n\bigl((w-a)^n - (z-a)^n\bigr)$，再对每一项除以
$w - z$（利用 $u^n - v^n$ 的因式分解）得到，对
$z, w \in D(a, r)$，
\[
g(z, w) =
\sum_{n\geq1}c_n\sum_{j=0}^{n-1}(w-a)^{\,j}(z-a)^{\,n-1-j},
\]
在对角线上同样成立（每个内和变为
$n(z-a)^{n-1}$，求和得 $f'(z)$）。当 $r$ 充分小时，该级数
在 $D(a,r)^2$ 上一致收敛
（$\abs{\text{项}} \leq n\abs{c_n}r^{n-1}$，
在收敛半径内可和），故其和连续。

令 $h(z) = \frac1{2\iu\pi}\int_\Gamma g(z, w)\,\dd w$ 在
$\Omega$ 上连续（由 $g$ 在紧集上一致连续）；
并且全纯——由莫雷拉定理
（\cref{thm:b3:holomorphic:weierstrassconv} 的判别准则）：对
三角形 $T \subseteq \Omega$，由富比尼定理得 $\int_{\partial
T}h = \frac1{2\iu\pi}\int_\Gamma\bigl(\int_{\partial T}g(z,
w)\dd z\bigr)\dd w = 0$，内层积分为零，因为
$z \mapsto g(z, w)$ 在 $\Omega$ 上全纯（在 $z = w$ 处
的奇点是可去的：$g$ 在该点连续，
在其他点全纯——即
\cref{thm:b3:holomorphic:formula} 证明中的延拓论证）。

在开集 $\Omega' = \{z \notin \operatorname{im}\Gamma
: \operatorname{Ind}_\Gamma(z) = 0\}$，令 $h_1(z) =
\frac1{2\iu\pi}\int_\Gamma\frac{f(w)}{w - z}\dd w$:
在 $\Omega'$ 上全纯（由莫雷拉定理或
对积分号求导）。对 $z \in \Omega\cap\Omega'$：
\[
h(z) = \frac1{2\iu\pi}\int_\Gamma\frac{f(w)}{w-z}\dd w
- f(z)\operatorname{Ind}_\Gamma(z) = h_1(z) .
\]
由假设 $\Omega\cup\Omega' = \C$（$w \notin \Omega
\Rightarrow \operatorname{Ind}_\Gamma(w) = 0$），故 $h$ 与
$h_1$ 粘合为整函数 $H$。由于
$\operatorname{im}\Gamma$ 补集的无界分支位于
$\Omega'$ 中，且 $h_1(z) \to 0$（当 $\abs z \to \infty$ 时，由 ML
估计），$H$ 有界且趋于 $0$：由刘维尔
定理（\cref{cor:b3:holomorphic:liouville}）得 $H \equiv 0$。于是
$h \equiv 0$ 在 $\Omega$ 上，即积分公式成立。
将其用于固定的 $a \in
\Omega\setminus\operatorname{im}\Gamma$，应用于 $\tilde f(w) =
(w - a)f(w)$，并取 $z = a$：
\[
\frac1{2\iu\pi}\int_\Gamma f(w)\dd w =
\frac1{2\iu\pi}\int_\Gamma \frac{\tilde f(w)}{w - a}\dd w =
\operatorname{Ind}_\Gamma(a)\,\tilde f(a) = 0 .
\]
\end{proof}

\section{洛朗级数与孤立奇点}

\begin{theorem}[洛朗展开]\label{thm:b3:residues:laurent}
设 $f$ 在圆环 $A = \{r < \abs{z - a} <
R\}$（$0 \leq r < R \leq \infty$）上全纯。则\index{洛朗级数}
\[
f(z) = \sum_{n\in\Z}c_n\,(z - a)^n
\qquad\text{ 在 } A,
\qquad
c_n = \frac1{2\iu\pi}\int_{C_\rho}\frac{f(w)}{(w -
a)^{n+1}}\,\dd w
\]
对任意 $r < \rho < R$（结果与 $\rho$ 无关），两个
半级数在紧子圆环上一致收敛。
该展开是唯一的。
\end{theorem}

\begin{proof}
固定 $r < \rho_1 < \abs{z - a} < \rho_2 < R$，并令 $\Gamma =
C_{\rho_2} - C_{\rho_1}$（外圆逆时针，内圆
顺时针）：$A$ 中的闭链，满足
$\operatorname{Ind}_\Gamma(w) = 0$，对所有 $w \notin A$
（小圆盘内的点：$1 - 1 = 0$；大圆外的
点：$0 - 0$）。由\cref{thm:b3:residues:globalcauchy}，
$\operatorname{Ind}_\Gamma(z) = 1 - 0 = 1$，于是
\[
f(z) = \frac1{2\iu\pi}\int_{C_{\rho_2}}\frac{f(w)}{w - z}\dd
w - \frac1{2\iu\pi}\int_{C_{\rho_1}}\frac{f(w)}{w - z}\dd w .
\]
对第一个核，仿照
\cref{thm:b3:holomorphic:analytic}，取 $\frac{z -
a}{w - a}$ 的幂（模 $< 1$）：非负部分
$\sum_{n\geq0}c_n(z-a)^n$。对第二项，改用另一种
方式展开：$\frac{-1}{w - z} = \frac1{(z-a)(1 - \frac{w - a}{z -
a})} = \sum_{m\geq0}\frac{(w-a)^m}{(z - a)^{m+1}}$，且
在 $C_{\rho_1}$ 上一致收敛：负幂部分
$\sum_{n\leq-1}c_n(z-a)^n$ 系数如上所述（指标
$n = -m-1$）。关于 $\rho$ 的无关性：
在 $C_{\rho}$ 与 $C_{\rho'}$ 上的系数积分相差
$\int_\Gamma$ 沿环绕数为零的闭链，对
$\frac{f(w)}{(w-a)^{n+1}}$（$A$ 中的全纯函数）的积分为零，
再次由\cref{thm:b3:residues:globalcauchy}。唯一性：
将 $\sum c_n(z-a)^n$ 乘以 $(z - a)^{-m-1}$，
再在 $C_\rho$ 上逐项积分（由一致收敛性），只有 $n = m$
的项保留下来。
\end{proof}

\begin{definition}\label{def:b3:residues:singularities}
若 $f$ 在去心圆盘 $D(a, R)\setminus
\{a\}$ 上全纯，用洛朗展开（取 $r = 0$）。三种互斥
情形：\index{孤立奇点}
\begin{itemize}
  \item 所有 $c_n = 0$（$n < 0$）：\emph{可去
        奇点}（非负幂级数将 $f$
        全纯延拓到 $a$）；
  \item $c_n \neq 0$ 仅对有限个、且至少一个 $n < 0$ 成立：
        为\emph{极点}，阶为 $m = -\min\{n : c_n \neq 0\}$；
        等价地，$f = g/(z-a)^m$，$g$ 全纯，$g(a)
        \neq 0$；等价地，$\abs{f(z)} \to \infty$，当 $z\to
        a$ 时；
  \item 满足 $c_n \neq 0$ 的负指标有无穷多个：
        \emph{本性奇点}。
\end{itemize}
\emph{留数}\index{留数}定义为
$\operatorname{Res}(f, a) = c_{-1}$。若函数在
$\Omega$ 中去掉某个极点集后全纯，则称该函数
为\emph{亚纯函数}\index{亚纯函数}（在 $\Omega$ 上）。
\end{definition}

\begin{theorem}[黎曼；卡索拉蒂--魏尔斯特拉斯]
\label{thm:b3:residues:riemanncw}
设 $f$ 在 $D(a,R)\setminus\{a\}$ 上全纯。
\begin{enumerate}
  \item （黎曼）若 $f$ 在 $a$ 附近\emph{有界}，
        则该奇点可去。\index{可去奇点}
  \item （卡索拉蒂--魏尔斯特拉斯）若 $a$ 为本性奇点，则
        $f\bigl(D(a,\varepsilon)\setminus\{a\}\bigr)$ 在
        $\C$ 中稠密，对每个 $\varepsilon$ 成立。
        \index{卡索拉蒂--魏尔斯特拉斯定理}
\end{enumerate}
\end{theorem}

\begin{proof}
（1）对 $n < 0$ 及 $\rho \to 0$：$\abs{c_n} \leq
M\rho^{-n-1}\cdot\rho\cdot\rho^{-\,n}\dots$ 由 ML 估计，
$C_\rho$: $\abs{c_n} \leq \frac{1}{2\pi}\,2\pi\rho\cdot
M\rho^{-(n+1)} = M\rho^{-n} \to 0$（因 $-n > 0$）：所有
负指标的系数全部为零。
（2）若某值 $b$ 未被逼近：$\abs{f - b} \geq
\delta$ 在 $a$ 附近，于是 $g = 1/(f - b)$ 全纯且
在 $a$ 附近有界，由（1）知可去，$g$ 延拓后的值为 $c$。
若 $c \neq 0$，$f = b + 1/g$ 在 $a$ 附近有界，故可去
--- 排除。若 $c = 0$，$g$ 的零点阶数有限，记为 $m$，
在 $a$ 处（见\cref{thm:b3:holomorphic:identity}；$g \not\equiv 0$），
则 $f = b + 1/g$ 有阶为 $m$ 的极点，同样排除。
\end{proof}

\section{留数定理}

\begin{theorem}[留数定理]\label{thm:b3:residues:residue}
设 $\Omega$ 为开集，$S \subseteq \Omega$ 为有限集，$f \in
\mathcal H(\Omega\setminus S)$，且 $\Gamma$ 是
$\Omega\setminus S$ 中的闭链，满足 $\operatorname{Ind}_\Gamma(w) = 0$
对每个 $w \notin \Omega$ 成立，则\index{留数定理}
\[
\frac{1}{2\iu\pi}\int_\Gamma f(z)\,\dd z
= \sum_{a\in S}\operatorname{Ind}_\Gamma(a)\,
\operatorname{Res}(f, a) .
\]
\end{theorem}

\begin{proof}
对每个 $a \in S$，令 $P_a(z) = \sum_{n\leq-1}c_n^{(a)}(z -
a)^n$ 为 $f$ 在 $a$ 处的主部：一个级数
在 $\C\setminus\{a\}$ 上收敛（它关于 $1/(z-a)$ 的
收敛半径为无穷：对所有充分小的 $\abs{z-a}$，洛朗尾部收敛，
因而作为 $(z-a)^{-1}$ 的幂级数处处收敛），并且
在该处全纯。于是 $g = f - \sum_{a\in S}P_a$ 在
$S$ 中每一点都是可去奇点（其洛朗
展开在 $a$ 处没有负幂部分：其余的 $P_{a'}$ 在
$a$ 处全纯），故 $g$ 可全纯延拓到
$\Omega$，由\cref{thm:b3:residues:globalcauchy}得
$\int_\Gamma g = 0$。下面对每个 $P_a$ 积分：
逐项积分（利用一致收敛性，在紧集
$\operatorname{im}\Gamma$ 上，它避开 $a$）：
\[
\frac1{2\iu\pi}\int_\Gamma(z - a)^n\,\dd z = 0 \ (n \leq -2:
\text{原函数 } \tfrac{(z-a)^{n+1}}{n+1}),
\qquad
\frac1{2\iu\pi}\int_\Gamma\frac{\dd z}{z - a} =
\operatorname{Ind}_\Gamma(a),
\]
于是 $\frac1{2\iu\pi}\int_\Gamma P_a =
c_{-1}^{(a)}\operatorname{Ind}_\Gamma(a)$。对 $a$ 求和。
\end{proof}

\begin{method}[留数的计算]\label{met:b3:residues:compute}
单极点：$\operatorname{Res}(f, a) = \lim_{z\to a}(z -
a)f(z)$；若 $f = g/h$ 且 $g(a) \neq 0$，$h(a) = 0$，$h'(a)
\neq 0$，则 $\operatorname{Res} = g(a)/h'(a)$。
$m$ 阶极点：$\operatorname{Res}(f, a) = \frac1{(m-1)!}\lim_{z\to
a}\bigl((z-a)^mf(z)\bigr)^{(m-1)}$。本性奇点：
展开并读出 $c_{-1}$（例如由已知级数得到）。始终需
验证围道实际围住了哪些极点，以及
各自的环绕数。
\end{method}

\begin{example}[四类经典积分]
\label{ex:b3:residues:integrals}
(a) \emph{$\R$ 上的有理函数}：对于 $\int_\R\frac{\dd x}{1 +
x^4}$，取上半平面内的大半圆周 $S_R$ 作围道，
其上被积函数为 $O(R^{-4})$，故
$\int_{S_R} \to 0$（ML），由留数定理及
极点 $\eu^{\iu\pi/4}, \eu^{3\iu\pi/4}$（均为单极点，留数
为 $\frac1{4z^3} = \frac{z}{4z^4} = -\frac z4$ 在极点处），于是得
\[
\int_\R\frac{\dd x}{1 + x^4}
= 2\iu\pi\Bigl(-\frac{\eu^{\iu\pi/4}}4 -
\frac{\eu^{3\iu\pi/4}}4\Bigr)
= \frac{\pi}{\sqrt2} .
\]
(b) \emph{傅里叶型}：对 $t \geq 0$，
$\int_\R\frac{\eu^{\iu tx}}{1 + x^2}\dd x =
2\iu\pi\operatorname{Res}\bigl(\tfrac{\eu^{\iu tz}}{1+z^2},
\iu\bigr) = 2\iu\pi\frac{\eu^{-t}}{2\iu} = \pi\eu^{-t}$ ---
取上半圆周可行，因为在其上 $\abs{\eu^{\iu tz}} =
\eu^{-t\operatorname{Im}z} \leq 1$；取实部得：
$\int_\R\frac{\cos(tx)}{1+x^2}\dd x = \pi\eu^{-\abs t}$，
由此证实\cref{exo:b3:lebesgue:10}中在本阶段不加证明地承认的公式。
(c) \emph{周期上的三角积分}：作代换 $z =
\eu^{\iu t}$, $\cos t = \frac{z + z^{-1}}2$, $\dd t =
\frac{\dd z}{\iu z}$: $\int_0^{2\pi}\frac{\dd t}{a + \cos t}$
（$a > 1$）时，该积分化为单位圆内的留数计算
（\cref{exo:b3:residues:2}）。
(d) \emph{级数}：将 $f$ 与 $\pi\cot(\pi z)$ 配对，其极点
为整数，留数均为 $1$：周末问题借此
求和 $\sum n^{-2}$ 与 $\sum n^{-4}$。
\end{example}

\begin{omfigure}
\begin{tikzpicture}[scale=1.15]
  \draw[->] (-2.5,0) -- (2.6,0) node[right, font=\small]
    {$\operatorname{Re}$};
  \draw[->] (0,-0.5) -- (0,2.5) node[above, font=\small]
    {$\operatorname{Im}$};
  \draw[omThm, very thick, ->] (-2,0) -- (0.3,0);
  \draw[omThm, very thick] (0.3,0) -- (2,0);
  \draw[omThm, very thick] (2,0) arc (0:90:2);
  \draw[omThm, very thick, ->] (2,0) arc (0:50:2);
  \draw[omThm, very thick] (0,2) arc (90:180:2);
  \node[omThm, font=\small] at (1.75,1.55) {$S_R$};
  \node[font=\small] at (2.2,-0.25) {$R$};
  \node[font=\small] at (-2.15,-0.25) {$-R$};
  \fill[omDef] (45:1.19) circle (1.6pt)
    node[above right, font=\small] {$\eu^{\iu\pi/4}$};
  \fill[omDef] (135:1.19) circle (1.6pt)
    node[above left, font=\small] {$\eu^{3\iu\pi/4}$};
  \fill[black!40] (-45:1.19) circle (1.6pt);
  \fill[black!40] (-135:1.19) circle (1.6pt);
\end{tikzpicture}

{\small 用于 $\int_\R\frac{\dd x}{1 +
x^4}$ 的半圆围道：当 $R \to \infty$ 时，圆弧贡献 $O(R^{-3})$，
留数定理计入被围住的两个极点（蓝色）。
下方两个极点（灰色）位于围道之外，指标为 $0$。}
\end{omfigure}

\section{辐角原理与儒歇定理}

\begin{theorem}[辐角原理]
\label{thm:b3:residues:argument}
设 $f$ 为 $\Omega$ 上的亚纯函数，零点为 $z_j$（阶数为
$m_j$），极点为 $p_k$（阶数为 $\mu_k$），又 $\gamma$ 为闭
路径，位于 $\Omega$ 内且避开所有这些点，并满足
$\operatorname{Ind}_\gamma = 0$ 于 $\Omega$ 之外，则
\index{辐角原理}
\[
\frac1{2\iu\pi}\int_\gamma\frac{f'(z)}{f(z)}\,\dd z
= \sum_j m_j\operatorname{Ind}_\gamma(z_j)
- \sum_k \mu_k\operatorname{Ind}_\gamma(p_k)
\]
（仅有限多项非零）。对于简单的正向
闭路，该积分计算内部零点个数减去
极点个数（计重数） --- 且等于
像路径 $f\circ\gamma$ 绕 $0$ 的环绕数。
\end{theorem}

\begin{proof}
在 $m$ 阶零点附近：$f = (z-a)^mg$，$g(a) \neq 0$，因此
$\frac{f'}f = \frac m{z - a} + \frac{g'}g$，其中第二项
在 $a$ 附近全纯，故为留数是 $m$ 的单极点。在
$\mu$ 阶极点附近：$f = (z-a)^{-\mu}g$，留数为
$-\mu$。其余处 $\frac{f'}f$ 全纯。（指标非零的零点与
极点位于被
$\gamma$ 所围的紧区域内；由恒等定理，它们在其中只有有限个，
因 $f \not\equiv 0$）。应用
\cref{thm:b3:residues:residue}。最后一点：
$\frac1{2\iu\pi}\int_\gamma\frac{f'}f =
\frac1{2\iu\pi}\int_{f\circ\gamma}\frac{\dd w}w =
\operatorname{Ind}_{f\circ\gamma}(0)$（作代换 $w =
f(\gamma(t))$）。
\end{proof}

\begin{theorem}[儒歇]\label{thm:b3:residues:rouche}
设 $f, g$ 为 $\Omega$ 上的全纯函数，$\gamma$ 为闭
路径，满足 $\operatorname{Ind}_\gamma \in \{0,1\}$，且在
$\Omega$ 之外为零（即简单闭路）。若\index{儒歇定理}
\[
\abs{g(z)} < \abs{f(z)} \qquad \text{在 }
\operatorname{im}\gamma,
\]
则 $f$ 与 $f + g$ 在
区域 $\{\operatorname{Ind}_\gamma =
1\}$ 内的零点个数（计重数）相同。
\end{theorem}

\begin{proof}
对 $t \in \intcc01$，$f_t = f + tg$ 在
$\operatorname{im}\gamma$ 上无零点（$\abs{f_t} \geq \abs f - \abs g >
0$），故
\[
N(t) = \frac1{2\iu\pi}\int_\gamma
\frac{f_t'(z)}{f_t(z)}\,\dd z
\]
良定义；它计数所围区域内的零点
（由\cref{thm:b3:residues:argument}，无极点）。$N$ 关于
$t$ 连续（被积函数联合连续，
分母一致有正下界 --- 控制
收敛定理），且取整数值，故为常数。$N(0) = N(1)$。
\end{proof}

\begin{corollary}[开映射定理]
\label{cor:b3:residues:openmapping}
连通开集上的非常值全纯函数为
开映射。特别地（再次）最大模原理
成立，全纯双射的
逆仍全纯。\index{开映射定理（全纯）}
\end{corollary}

\begin{proof}
设 $f(a) = b$；$f - b$ 有有限阶 $m
\geq 1$ 的零点，位于 $a$ 处（由恒等定理，$f \not\equiv b$）。取
$r$，使 $f - b$ 在 $\bar D(a, r)\setminus\{a\}$ 上
无零点（零点孤立），令 $\delta = \min_{\abs{z - a} =
r}\abs{f(z) - b} > 0$。对 $\abs{w - b} < \delta$，在
圆周上，$\abs{(b - w)} < \delta \leq \abs{f - b}$，
故由儒歇定理（比较 $f - b$ 与常数 $b - w$）可知 $f - w$ 恰有
$m$ 个零点，位于 $D(a, r)$ 内：每个这样的 $w$ 均被取到
--- $f(D(a,r)) \supseteq D(b, \delta)$，故为开映射。最大模
原理：$\abs f$ 在内点处取极大值是不可能的，
因为对非常值 $f$，其在 $f(a)$ 附近的像含有
模更大的点。逆：全纯双射 $f$ 为
开映射，故 $f^{-1}$ 连续；$f - f(a)$ 在
$a$ 处的零点是单的（若 $m \geq 2$，则附近的值将有 $m$ 个原像
--- 且互异，因为 $f'$ 仅在
孤立点处为零，故在 $a$ 附近，$m$ 个
$f - w$ 的零点对一般的小 $w$ 均为单的且互异，与
单射性矛盾）；于是 $f'(a) \neq 0$，且 $f^{-1}$ 的差
商收敛：$\bigl(f^{-1}\bigr)'(b) =
1/f'(a)$。
\end{proof}

\section{练习}

\begin{exercise}[$\star$]\label{exo:b3:residues:1}
判断 $0$ 处奇点的类型，并求留数：
\[
\frac{\sin z}{z},\qquad
\frac{\eu^z - 1}{z^2},\qquad
\frac{1}{z(z-1)^2},\qquad
\frac{\cos z}{z^3},\qquad
\eu^{1/z},\qquad
\frac1{\sin z} .
\]
另求第三个函数在 $z = 1$ 处的留数，以及最后一个函数
在 $z = \pi$ 处的留数。
\end{exercise}

\begin{exercise}[$\star$]\label{exo:b3:residues:2}
设 $a > 1$，借助 $z = \eu^{\iu t}$ 计算：
\[
\int_0^{2\pi}\frac{\dd t}{a + \cos t}
= \frac{2\pi}{\sqrt{a^2 - 1}} .
\]
检验 $a \to 1^+$ 与 $a \to \infty$ 时的极限行为。
\end{exercise}

\begin{exercise}[$\star\star$]\label{exo:b3:residues:3}
利用半圆形积分路径计算，并论证弧上的估计：
\[
\int_\R\frac{x^2}{1 + x^6}\,\dd x = \frac\pi3,
\qquad
\int_\R\frac{\dd x}{(1 + x^2)^{2}} = \frac\pi2
\quad\text{（重新推导 \cref{exo:b3:fouriertransform:3}）}.
\]
\end{exercise}

\begin{exercise}[$\star\star$]\label{exo:b3:residues:4}
对 $t \geq 0$ 与 $a > 0$，证明：
\[
\int_\R\frac{\cos(tx)}{x^2 + a^2}\,\dd x =
\frac{\pi}{a}\,\eu^{-at},
\]
并由反演求函数 $x \mapsto \eu^{-a\abs
x}$ 的傅里叶变换 --- 并与
\cref{exo:b3:fouriertransform:1} 比较。
\end{exercise}

\begin{exercise}[$\star\star$]\label{exo:b3:residues:5}
将 $f(z) = \dfrac1{(z-1)(z-2)}$ 分别在
三个区域 $\abs z < 1$、$1 < \abs z < 2$、$\abs z >
2$ 内展开为洛朗级数。三种展开为何不同？说明为何
第二、第三个展开式中 $z^{-1}$ 的系数
\emph{不是} $f$ 在 $0$ 处的留数（$f$ 在该处无奇点
），并求 $f$ 在 $1$ 与
$2$ 处的实际留数。
\end{exercise}

\begin{exercise}[$\star\star$]\label{exo:b3:residues:6}
(a) 证明 $\eu^{1/z}$ 在 $0$ 处有本性奇点，
并直接验证卡索拉蒂--魏尔斯特拉斯定理：解方程 $\eu^{1/z} =
w$ 的显式解（对任意 $w \neq 0$），给出
任意接近 $0$ 的解。
(b) 证明 $\abs{\eu^{1/z}}$ 在每个
$0$ 的去心邻域上无界，但 $\eu^{1/z}$ 无极点：
哪个极限失效？
\end{exercise}

\begin{exercise}[$\star\star$]\label{exo:b3:residues:7}
用儒歇定理计数：
(a) 求 $z^7 - 4z^3 + z - 1$ 在 $\abs z < 1$ 内的零点个数；
(b) 求 $z^4 + 5z + 1$ 在 $\abs z < 1$ 内及在 $1 <
\abs z < 2$ 内的零点个数；
(c) 重新证明达朗贝尔–高斯定理：首一的 $n$ 次多项式
在某个足够大的圆盘内有 $n$ 个零点 \emph{（与
$z^n$ 比较）}。
\end{exercise}

\begin{exercise}[$\star\star\star$]\label{exo:b3:residues:8}
（赫尔维茨）设 $f_n \to f$ 在紧子集上一致收敛，$f_n \in
\mathcal H(\Omega)$，$\Omega$ 连通，$f \not\equiv 0$。
(a) 证明：若所有 $f_n$ 都无零点，则 $f$ 也无零点。
\emph{（若 $f(a) = 0$：对 $a$ 附近的小圆周应用辐角原理，
并由 \cref{thm:b3:holomorphic:weierstrassconv} 对
$\int f_n'/f_n$ 取极限。）}
(b) 证明：若所有 $f_n$ 都是单射，则 $f$ 为单射
或常值函数。\emph{（对 $z \mapsto f_n(z) - f_n(w)$ 在
$\Omega\setminus\{w\}$ 上应用 (a)。）}
\end{exercise}

\begin{exercise}[$\star\star\star$]\label{exo:b3:residues:9}
对 $n \geq 2$，将 $\frac1{1 + z^n}$ 沿
扇形 $\{0 \leq \arg z \leq \frac{2\pi}n,\ \abs z
\leq R\}$ 的边界积分，并推得
\[
\int_0^{+\infty}\frac{\dd x}{1 + x^n} =
\frac{\pi}{n\,\sin(\pi/n)} .
\]
验证 $n = 2$ 时结果与 $\arctan$ 一致，并考虑极限 $n \to
\infty$ 的情形。
\end{exercise}

\begin{exercise}[$\star\star$]\label{exo:b3:residues:10}
设 $f$ 为有理函数，满足 $\deg(\text{分母})
\geq \deg(\text{分子}) + 2$，证明 $f$ 的
\emph{全部}留数之和为零\emph{（沿
越来越大的圆周积分）}。利用此结论，
重新计算 $\frac1{z(z-1)(z-2)}$ 的
部分分式分解，不借助线性代数。
\end{exercise}

\begin{exercise}[$\star\star\star$]\label{exo:b3:residues:11}
（钥匙孔围道：欧拉反射积分）对 $0 < a < 1$，
计算
\[
I(a) = \int_0^{\infty}\frac{x^{a-1}}{1 + x}\,\dd x =
\frac{\pi}{\sin(\pi a)}
\]
通过对 $f(z) = \frac{z^{a-1}}{1+z} =
\frac{\eu^{(a-1)\log z}}{1 + z}$（对数沿
$\R_+$ 割开，$\arg z \in \intoo0{2\pi}$），在钥匙孔
围道上积分：沿割线上侧从 $\varepsilon$ 到
$R$，绕 $C_R$，再沿割线下侧返回，绕
$C_\varepsilon$。说明：两段直线部分相差
因子 $\eu^{2\iu\pi(a-1)}$，圆弧部分的贡献趋于零
（$R^{a-1}\cdot R \to 0$ 与
$\varepsilon^{a-1}\cdot\varepsilon \to 0$），唯一极点 $z = -1$ 的
留数为 $\eu^{\iu\pi(a - 1)}$。由此还可推得
$\Gamma(a)\Gamma(1 - a) = \frac\pi{\sin\pi a}$ \emph{（先写
$\Gamma(a)\Gamma(1-a) = B(a, 1-a)$，由
\cref{pb:b3:lebesgue:1}，并作代换 $t =
\frac{x}{1+x}$）}。
\end{exercise}

\begin{exercise}[$\star\star$]\label{exo:b3:residues:12}
（用辐角原理计数零点，并作数值计算）
设 $P(z) = z^4 + 8z + 1$。
(a) 单位圆盘内有多少个零点？\emph{（用儒歇定理，
与 $8z + 1$ 比较。）}
(b) 在圆环 $1 < \abs z < 3$ 内有多少个零点？
\emph{（用儒歇定理，与 $z^4$ 在 $\abs z = 3$ 上比较。）}进一步细化：
证明每个零点的模都 $< 2.1$。
(c) 右半平面内有多少个零点？\emph{（先在
$\abs z = 2$ 上计数；再考察虚轴的
像：$P(\iu t) = t^4 + 1 + 8\iu t$ 的实部处处为正，
因此虚轴上没有零点，且沿虚轴的辐角
变化可以算出 --- 最后以一个大的
半圆盘收尾。）}
\end{exercise}

\section{问题：\texorpdfstring{$\zeta(2k)$}{zeta(2k)}
借助余切函数}

\begin{problem}[{周末问题 --- 用留数求 $\sum n^{-2k}$
的和}]\label{pb:b3:residues:1}
留数定理可用于求级数之和：将有理函数与
$\pi\cot(\pi z)$ 配对，其极点位于整数处，使
$\sum_{n}f(n)$ 化为留数计数。我们证明该方法，
并计算 $\zeta(2) = \frac{\pi^2}{6}$ 与 $\zeta(4) =
\frac{\pi^4}{90}$ --- 这两个值在本科二年级用傅里叶级数求得，
又在\cref{ch:b3:spectral}中由算子迹求得，现以
围道积分重新求得。

\medskip
\noindent\textbf{第一部分 --- 余切核。}
\begin{enumerate}
  \item 证明 $\pi\cot(\pi z)$ 在 $\C$ 上亚纯，且
        恰在 $z = n \in \Z$ 处有单极点，每个极点的
        留数为 $1$ \emph{（计算
        $\lim_{z\to n}(z-n)\pi\cot\pi z$）}。
  \item 求 $0$ 处洛朗展开的开头几项：
        \[
        \pi\cot(\pi z) = \frac1z - \frac{\pi^2}{3}\,z -
        \frac{\pi^4}{45}\,z^3 + O(z^5) ,
        \]
        用 $\cos$ 的幂级数除以
        $\sin$ 的幂级数（说明除法可行：$\frac{\sin\pi z}{\pi z}$
        在 $0$ 附近全纯且不为零，故
        其倒数也全纯；比较各项系数，
        直到 $3$ 次项为止）。
  \item 设 $C_N$ 为正方形的边界，其顶点为
        $(\pm1\pm\iu)(N + \frac12)$。证明
        $\abs{\cot(\pi z)} \leq 2$ 在 $C_N$ 上对一切 $N \geq
        1$ 成立。\emph{（竖直边上，$\cot(\pi(\pm(N +
        \frac12) + \iu y)) = \mp\tan(\iu\pi y)$，其模为
        $\abs{\tanh(\pi y)} \leq 1$；水平边上
        $\abs y = N + \frac12$，估计
        $\abs{\cot(\pi(x\pm\iu y))} \leq \coth(\pi y) \leq
        \coth(\pi/2) < 1.1$。）}
\end{enumerate}

\medskip
\noindent\textbf{第 II 部分 --- 求和定理。}
\begin{enumerate}[resume]
  \item 设 $f$ 为有理函数，在整数处全纯，
        且 $\deg(\text{分母}) \geq \deg(\text{分子}) + 2$。
        利用 $C_N$ 上的留数定理，
        及第 3 题中的估计，证明：
        \[
        \lim_{N\to\infty}\ \sum_{n = -N}^{N} f(n)
        = -\sum_{p\ \text{极点}\ f}
        \operatorname{Res}\bigl(\pi\cot(\pi z)f(z),\,p\bigr).
        \]
  \item 论证在何处需要次数条件？
        通过例子说明（取 $f(z) = 1/(z + \frac12)$），
        当衰减较慢时，对称极限
        仍可能存在，而双边级数发散，
        --- 并说明公式此时计算的是
        \emph{主值}。
\end{enumerate}

\medskip
\noindent\textbf{第 III 部分 --- 取值。}
\begin{enumerate}[resume]
  \item 将方法应用于 $f(z) = 1/z^2$：此时 $f$ 的
        极点\emph{恰位于}整数处，故直接进行论证
        --- 对 $g(z) = \frac{\pi\cot(\pi
        z)}{z^2}$ 沿 $C_N$ 积分，证明该积分 $\to 0$，
        并由第 2 题计算 $\operatorname{Res}(g, 0)$。
        于是得到：
        \[
        2\sum_{n\geq1}\frac1{n^2} = \frac{\pi^2}{3},
        \qquad \zeta(2) = \frac{\pi^2}6 .
        \]
  \item 对 $g(z) = \frac{\pi\cot(\pi z)}{z^4}$ 作同样处理：
        计算 $\operatorname{Res}(g, 0)$，并推出
        $\zeta(4) = \frac{\pi^4}{90}$。
  \item 说明一般规律：对每个 $k \geq 1$，
        $\zeta(2k)$ 为 $-\frac12$ 倍的
        $z^{2k-1}$ 项系数，该系数取自 $\pi\cot(\pi
        z)$ 在 $0$ 处的洛朗展开 --- 即 $\pi^{2k}$ 的一个有理倍数。
        将第 2 题中的除法
        再推进一步，计算 $\zeta(6)$。该方法对
        $\zeta(3)$ 有何说法？为何它什么也说不出？
\end{enumerate}

\medskip
\noindent\textbf{第 IV 部分 --- 余切函数的部分分式
展开。}
\begin{enumerate}[resume]
  \item 固定 $w \in \C\setminus\Z$，并将
        第 II 部分的方法应用于 $f(z) = \dfrac{1}{(z - w)(z + w)}$ ---
        注意此时 $\pi\cot(\pi z)f(z)$ 另外还有
        在 $\pm w$ 处的单极点，其留数须计入
        计数。由此推出部分分式展开
        \[
        \pi\cot(\pi w) = \frac1w +
        \sum_{n\geq1}\frac{2w}{w^2 - n^2},
        \]
        该级数在
        $\C\setminus\Z$ 的紧子集上正规收敛。
  \item 由此展开式，将每一项展开为
        $w$ 的幂级数（需说明交换求和次序的依据），得到
        与第 2 题中\emph{相同}的洛朗系数 ---
        至此闭合：欧拉公式
        $\sum\frac1{n^2} = \frac{\pi^2}6$ 即是
        余切函数两种面貌中 $w$ 项的系数。将其与
        本科二年级的傅里叶级数证明，以及迹证明
        （\cref{pb:b3:spectral:1}）：三种理论，
        同一个数。
\end{enumerate}

\medskip
\noindent\textbf{第 V 部分 --- 正弦函数的欧拉乘积。}
第 9 题中的展开式，是某个无穷乘积的
对数导数。现在我们严格证明欧拉于 1734 年给出的
因式分解。
\begin{enumerate}[resume]
  \item 对 $N \geq 1$，记 $P_N(z) =
        z\prod_{n=1}^{N}\bigl(1 - \frac{z^2}{n^2}\bigr)$。
        证明 $P_N$ 在每个圆盘
        $\bar D(0, R)$ 上一致收敛于整函数 $P$，其零点
        恰为全体整数，且均为单零点。\emph{（对 $n
        \geq 2R$ 时，把因子写成 $\exp\log(1 -
        z^2/n^2)$，取主对数（见
        \cref{exo:b3:holomorphic:3}），估计 $\abs{\log(1+u)}
        \leq 2\abs u$（当 $\abs u \leq \frac12$ 时），由
        幂级数，并对正规收敛的
        对数和取指数；其余有限个因子
        为多项式。由
        \cref{thm:b3:holomorphic:weierstrassconv}即得结论。）}
  \item 证明在 $\C\setminus\Z$ 上，
        \[
        \frac{P'(z)}{P(z)} = \frac1z +
        \sum_{n\geq1}\frac{2z}{z^2 - n^2} = \pi\cot(\pi z)
        \]
        \emph{（对有限乘积求导，再取极限，
        利用
        \cref{thm:b3:holomorphic:weierstrassconv}及
        $P$ 在 $\Z$ 之外无零点，并引用第
        9 题）}。
  \item 证明 $Q = \sin(\pi z)/P(z)$ 可延拓为
        无零点的整函数，且 $Q' = 0$，
        由此得出\emph{欧拉乘积}：
        \[
        \sin(\pi z) = \pi z\prod_{n\geq1}
        \Bigl(1 - \frac{z^2}{n^2}\Bigr)
        \qquad (z \in \C) .
        \]
  \item （沃利斯，1655 年）在 $z = \frac12$ 处求值：
        \[
        \frac\pi2 = \prod_{n\geq1}\frac{4n^2}{4n^2 - 1}
        = \lim_{N\to\infty}
        \frac{2\cdot2\cdot4\cdot4\cdots(2N)(2N)}
        {1\cdot3\cdot3\cdot5\cdots(2N-1)(2N+1)} .
        \]
  \item 对 $\abs z < 1$，将乘积的对数展开为
        二重级数（需说明重排的依据），
        由此推出 $\zeta(2) =
        \frac{\pi^2}6$，即比较 $z^3$ 的系数，
        在 $\sin(\pi z) = \pi z - \frac{\pi^3}6z^3 + \cdots$ 的展开式中，
        --- 这正是欧拉数在乘积形式下的面貌。
\end{enumerate}

\medskip
\noindent\textbf{第六部分 --- 同族核。}余切核
有若干同族核；每个核各自为一族级数定价。
\begin{enumerate}[resume]
  \item 对第 9 题的展开式逐项求导
        （依据
        \cref{thm:b3:holomorphic:weierstrassconv}），得到
        在 $\C\setminus\Z$ 的紧子集上正规收敛：
        \[
        \frac{\pi^2}{\sin^2(\pi z)} =
        \sum_{n\in\Z}\frac1{(z - n)^2} .
        \]
  \item 令 $z = \frac12$，得：
        $\sum_{m\geq0}\frac1{(2m+1)^2} = \frac{\pi^2}8$;
        按奇偶性拆分整数，再次恢复
        $\zeta(2)$。
  \item 验证倍角恒等式 $\tan\theta =
        \cot\theta - 2\cot(2\theta)$，并推出
        \[
        \pi\tan(\pi z) =
        \sum_{m\geq0}\frac{8z}{(2m+1)^2 - 4z^2} ,
        \]
        在避开 $\frac12 + \Z$ 的紧集上正规收敛。
  \item 在 $0$ 附近展开（$\abs z < \frac12$；
        再次使用富比尼定理），其中 $\lambda(s) =
        \sum_{m\geq0}(2m+1)^{-s}$,
        \[
        \pi\tan(\pi z) =
        \sum_{k\geq0}8\cdot4^k\,\lambda(2k+2)\,z^{2k+1} ;
        \]
        与 $\tan u = u + \frac{u^3}3 + O(u^5)$ 比较，以
        恢复 $\lambda(2) = \frac{\pi^2}8$，并求得
        $\lambda(4) = \frac{\pi^4}{96}$，再
        验证 $\zeta(4) = \frac{\pi^4}{90}$：由 $\lambda(4) = (1 -
        2^{-4})\,\zeta(4)$ 可得。
  \item 验证 $\frac1{\sin\theta} = \cot\frac\theta2 -
        \cot\theta$，并推出
        \[
        \frac{\pi}{\sin(\pi z)} = \frac1z +
        \sum_{n\geq1}(-1)^n\,\frac{2z}{z^2 - n^2} .
        \]
        借助
        $\pi/\sin(\pi z)$ 在整数点处的留数检验符号。
  \item 读出 $z$ 的系数：$\eta(2) =
        \sum_{n\geq1}\frac{(-1)^{n-1}}{n^2} =
        \frac{\pi^2}{12}$，并验证
        一致性 $\eta(2) = (1 - 2^{1-2})\,\zeta(2)$。
  \item （压轴）将第 20 题的展开式在 $z =
        \frac12$ 处求值，并推出 \emph{莱布尼茨公式}
        \[
        \frac\pi4 = 1 - \frac13 + \frac15 - \frac17 +
        \cdots
        \]
        以一段简短文字作结：每个核对应一种算术
        --- 哪一个核为哪一族级数定价，
        以及为何它们在结构上都对
        $\zeta(3)$ 视而不见。

\end{enumerate}
\medskip
\noindent\textbf{第七部分 --- 完整价目表：伯努利
数。}
\begin{enumerate}[resume]
  \item 将 $\pi
        z\cot(\pi z)$ 的部分分式展开与伯努利数的母函数相结合
        （$\frac{w}{\eu^w - 1} =
        \sum_n\frac{B_n}{n!}w^n$,
        见\cref{pb:b3:holomorphic:1}，第六部分）：由
        \[
        \pi z\cot(\pi z) = \iu\pi z +
        \frac{2\iu\pi z}{\eu^{2\iu\pi z} - 1}
        \]
        （先证明此恒等式），推出闭形式
        \[
        \zeta(2k) = (-1)^{k+1}\,
        \frac{(2\pi)^{2k}\,B_{2k}}{2\,(2k)!}
        \qquad (k \geq 1).
        \]
  \item 对照 $B_2 = \frac16$，$B_4 =
        -\frac1{30}$，$B_6 = \frac1{42}$，检验公式，恢复 $\zeta(2) =
        \frac{\pi^2}6$，$\zeta(4) = \frac{\pi^4}{90}$，并
        计算 $\zeta(6) = \frac{\pi^6}{945}$。
  \item （欧拉递推）将
        $\bigl(z\cot z\bigr)' = \cot z - z(1 + \cot^2z)$ 两边展开 ---
        或直接平方余切级数 --- 以证明
        \[
        \Bigl(k + \frac12\Bigr)\zeta(2k) =
        \sum_{j=1}^{k-1}\zeta(2j)\,\zeta(2k - 2j)
        \qquad (k \geq 2),
        \]
        并检验 $\zeta(4)$ 可由 $\zeta(2)$ 算得，
        $\zeta(6)$ 可由 $\zeta(2), \zeta(4)$ 算得 --- 全部偶数
        阶的 zeta 值都只由单一起点 $\frac{\pi^2}6$ 得到，
        无需任何新的积分。
\end{enumerate}
\end{problem}
